\documentclass[11pt]{article}

\usepackage[T1]{fontenc}
\usepackage{lmodern}
\usepackage[margin=0.92in]{geometry}
\usepackage{microtype}
\usepackage{amsmath,amssymb,amsthm,mathtools}
\usepackage{booktabs,longtable,array,tabularx}
\usepackage{xcolor}
\usepackage{enumitem}
\usepackage{fancyhdr}
\usepackage{xurl}
\usepackage{hyperref}
\usepackage[nameinlink,noabbrev]{cleveref}
\usepackage{listings}
\usepackage{bm}
\usepackage{needspace}
\usepackage{graphicx}
\usepackage{float}

\graphicspath{{figures/}}

% Keep hyperref's explicit document metadata while omitting volatile
% engine-generated timestamps, trailer identifiers, and PTEX banner fields.
\ifdefined\pdfinfoomitdate
  \pdfinfoomitdate=1
\fi
\ifdefined\pdftrailerid
  \pdftrailerid{}
\fi
\ifdefined\pdfsuppressptexinfo
  \pdfsuppressptexinfo=15
\fi

\allowdisplaybreaks
\Urlmuskip=0mu plus 1mu

\hypersetup{
  pdftitle={Projector-Resolved Seed Calculus III: Typed Edge Operators, Connection Defects, and Obstruction Packets},
  pdfauthor={Komorebi P.},
  pdfsubject={Typed edge operators, connection defects, holonomy, and obstruction packets},
  colorlinks=true,
  linkcolor=black,
  citecolor=black,
  urlcolor=black
}

\pagestyle{fancy}
\fancyhf{}
\fancyhead[L]{Projector-Resolved Seed Calculus III}
\fancyhead[R]{Komorebi P.}
\fancyfoot[C]{\thepage}
\setlength{\headheight}{14pt}

\newtheorem{theorem}{Theorem}[section]
\newtheorem{proposition}[theorem]{Proposition}
\newtheorem{lemma}[theorem]{Lemma}
\newtheorem{corollary}[theorem]{Corollary}
\theoremstyle{definition}
\newtheorem{definition}[theorem]{Definition}
\newtheorem{example}[theorem]{Example}
\newtheorem{construction}[theorem]{Construction}
\theoremstyle{remark}
\newtheorem{remark}[theorem]{Remark}
\newtheorem{boundary}[theorem]{Boundary}
\newtheorem{roadmap}[theorem]{Roadmap item}

\newcommand{\bbk}{\Bbbk}
\newcommand{\N}{\mathbb N}
\newcommand{\cA}{\mathcal A}
\newcommand{\cC}{\mathcal C}
\newcommand{\cE}{\mathcal E}
\newcommand{\cG}{\mathcal G}
\newcommand{\cH}{\mathcal H}
\newcommand{\cO}{\mathcal O}
\newcommand{\cP}{\mathcal P}
\newcommand{\cQ}{\mathcal Q}
\newcommand{\cS}{\mathcal S}
\newcommand{\cT}{\mathcal T}
\newcommand{\cW}{\mathcal W}
\newcommand{\End}{\operatorname{End}}
\newcommand{\Hom}{\operatorname{Hom}}
\newcommand{\im}{\operatorname{im}}
\newcommand{\id}{\mathrm{id}}
\newcommand{\Ob}{\operatorname{Ob}}
\newcommand{\Hol}{\operatorname{Hol}}
\newcommand{\Coord}{\operatorname{Coord}}
\newcommand{\ev}{\operatorname{ev}}
\newcommand{\src}{\operatorname{src}}
\newcommand{\tgt}{\operatorname{tgt}}
\newcommand{\Walk}{\operatorname{Walk}}
\newcommand{\supp}{\operatorname{supp}}
\newcommand{\op}{\mathrm{op}}
\newcommand{\leanname}[1]{\nolinkurl{#1}}
\newcolumntype{L}[1]{>{\raggedright\arraybackslash}p{#1}}
\newcolumntype{Y}{>{\raggedright\arraybackslash}X}

\definecolor{leanblue}{RGB}{20,72,135}
\definecolor{roadmapgray}{RGB}{86,86,86}

\newcommand{\statusleanthree}[1]{%
  \par\smallskip\noindent
  {\color{leanblue}\textbf{Formal status: Lean-backed in Paper III.}}\par
  \begingroup\small\raggedright\leanname{#1}\par\endgroup\smallskip}
\newcommand{\statusleanprior}[1]{%
  \par\smallskip\noindent
  {\color{leanblue}\textbf{Formal status: inherited Lean-backed result.}}\par
  \begingroup\small\raggedright\leanname{#1}\par\endgroup\smallskip}
\newcommand{\statusroadmap}[1]{%
  \par\smallskip\noindent
  {\color{roadmapgray}\textbf{Formal status: roadmap only.}}
  \ #1\par\smallskip}

\lstset{
  basicstyle=\ttfamily\footnotesize,
  breaklines=true,
  frame=single,
  showstringspaces=false,
  columns=fullflexible,
  keepspaces=true
}

\title{\textbf{Projector-Resolved Seed Calculus III}\\[0.45em]
\large Typed Edge Operators, Connection Defects, and Obstruction Packets}
\author{Komorebi P.}
\date{Version 1.0.1 --- September 3, 2026}

\begin{document}
\maketitle

\begin{abstract}
This paper separates three structures that are easy to conflate.
First, a finite packet graph supports typed elementary operators: an edge from
packet $i$ to packet $j$ contributes only to the $(j,i)$ matrix block.
Noncomposable products therefore vanish, and powers of the total transfer
operator expand into ordered operator words whose surviving terms are
directed walks.  Lean checks the recursive, explicit-word, and walk-filtered
forms, the endpoint coefficient formula, and fixed-length and orbit-support
reachability obstructions.
Second, exact packet coordinates on one ambient module carry a canonical
transport
\[
J_{b\leftarrow a}=\cA_b\cS_a.
\]
Its frame-level identity, cocycle, round-trip, and triangle-loop laws are
checked in Paper III; the released formalization also checks composition and
naturality for arbitrary functions between exact coordinate systems.  Linear
packet-coordinate counterparts remain available from the frozen contract of
Paper II.
Consequently canonical coordinate transport has no holonomy.  Third, a
connection $\Gamma_{b\leftarrow a}$ is additional, noncanonical transport
data.  Lean checks both a minimal self-normalized API and a heterogeneous
reversible API, including the additive defect--holonomy identity, finite-path
collapse under flatness, gauge covariance, and gauge equivalence of flat pair
connections to canonical transport.  It also checks exact finite block
reconstruction of triangle obstructions, the four equivalent vanishing
criteria, packet-supported witnesses, and packet-preserving block-gauge
covariance.  None of these results identifies $\Gamma$ with $J$ by definition.
Finally, Lean supplies finite unary--binary operation-tree
syntax, arbitrary semantics, locality, and candidate-specific obstruction
certificates.  Colored trees, multilinear packet expansions, and nonlinear
Fa\`a di Bruno defects remain roadmap items.  No general nonlinear chain rule,
analytic convergence theorem, or physical interpretation is claimed.
\end{abstract}

\begin{center}
\fbox{\begin{minipage}{0.93\textwidth}
\textbf{Final verification boundary.}
Version 1.0.1 is theorem-complete for its explicitly scoped claims.  All
thirty theorem-like environments are formally backed: twenty-six by compiled
Paper III declarations and four by the frozen Papers I--II contracts.  The
Paper III scope is checked with Lean 4 v4.34.0-rc2 against Mathlib revision
\leanname{aeb72f1e9981b6d35b426b16dca7759be3e7ba06}
\cite{lean,mathlib} from
\path{EdgeOperators.lean}, \path{EdgeExtensions.lean},
\path{Transport.lean}, \path{ReversibleConnections.lean},
\path{CoordinateNaturality.lean}, \path{HolonomyBlocks.lean}, and
\path{NonlinearObstructions.lean}; the declaration-level ledger is given in
\cref{sec:ledger}.  Roadmap items are explicit exclusions, not completion
claims.  Definitions do not count as theorems.  In particular, the words
\emph{connection}, \emph{curvature}, and \emph{holonomy} refer only to the
algebraic APIs defined here.  Formal verification establishes conformance to
the stated specifications; it does not by itself establish novelty or
priority.
\end{minipage}}
\end{center}

\clearpage
\tableofcontents
\clearpage

\section{Purpose and status discipline}
\label{sec:purpose}

Paper I fixed typed seeds, operator grammars, exact projector packets,
generated hulls, and reconstruction \cite{paperI}.  Paper II added coherent
refinement coordinates, weighted transfer, arithmetic separation, and linear
defect identities \cite{paperII}.  The present paper begins the next layer:
explicit edge words, genuinely optional connection data, and packet-local
obstructions.

The central distinction is categorical rather than terminological.
Canonical coordinate transport compares two exact descriptions of the
\emph{same} ambient object and is forced by analysis and synthesis.  A
connection chooses transports that may carry extra modeling, discretization,
or gauge information.  The former obeys a cocycle theorem; failure of the
latter to obey a cocycle is meaningful data.  Treating both maps as one
symbol would make holonomy either tautologically trivial or logically
unfounded.

\subsection{Statement labels}

Every substantial theorem, lemma, corollary, proposition, definition/API, and
roadmap item carries status evidence drawn from one or more of three labels.
A statement
whose separate clauses use both a Paper III result and an inherited bridge
carries both applicable labels.  Illustrative examples are exempt.
\begin{enumerate}[label=\textbf{S\arabic*.},leftmargin=3.2em]
\item \textbf{Lean-backed in Paper III}: the displayed scope is no stronger
than the named compiled declaration in the released Paper III module.
\item \textbf{Inherited Lean-backed}: the displayed result is supplied by the
frozen formal contract of Paper I or II.
\item \textbf{Roadmap only}: the text states a future interface or theorem
target and is excluded from the completion denominator.
\end{enumerate}

\begin{boundary}[Final theorem denominator]
The completion denominator is the thirty theorem-like environments in the
paper: twenty-six resolve to compiled Paper III declarations and four resolve
to the frozen Papers I--II contracts.  The authoritative theorem map names the
exact declarations used by each statement.  The kernel audit checks those
targets and supporting theorems, while a fail-closed gate scans the complete
compiled \leanname{Komorebi.PaperIII} namespace.  Roadmap items are excluded
from the denominator and remain visibly marked.
\end{boundary}

\section{Finite typed packet graphs}
\label{sec:edge-operators}

\subsection{Orientation and matrix convention}

Let $V$ be a finite packet type, $E$ a finite edge type, and $R$ a semiring.
A directed multigraph is a pair of endpoint functions
\[
  \src,\tgt:E\longrightarrow V.
\]
Parallel edges remain distinct elements of $E$.  Rows are targets and columns
are sources.  Thus an edge $e:i\to j$ is represented by a matrix supported at
row $j$, column $i$.

\begin{definition}[Packet graph and scalar edge system]
\label{def:edge-system}
For $a,b\in V$ and $r\in R$, define
\[
  \operatorname{packetMatrix}(a,b;r)_{yx}
  =
  \begin{cases}
  r,&y=b\text{ and }x=a,\\
  0,&\text{otherwise}.
  \end{cases}
\]
An edge system is a coefficient function $c:E\to R$.  Its elementary edge
operator and total transfer operator are
\[
  L_e=\operatorname{packetMatrix}(\src(e),\tgt(e);c(e)),
  \qquad
  T=\sum_{e\in E}L_e.
\]
\end{definition}
\statusleanthree{Komorebi.PaperIII.EdgeOperators.PacketGraph; Komorebi.PaperIII.EdgeOperators.PacketGraph.packetMatrix; Komorebi.PaperIII.EdgeOperators.PacketGraph.EdgeSystem; Komorebi.PaperIII.EdgeOperators.PacketGraph.EdgeSystem.edgeOp; Komorebi.PaperIII.EdgeOperators.PacketGraph.EdgeSystem.transferOp}

The convention is execution ordered.  In $L_fL_e$, the rightmost operator
$L_e$ acts first.  For a written word $(e_1,\ldots,e_n)$, $e_1$ acts first
and the operator product is $L_{e_n}\cdots L_{e_1}$.

\begin{lemma}[Declared support]
\label{lem:edge-support}
For every edge $e$, $(L_e)_{yx}=0$ unless
$y=\tgt(e)$ and $x=\src(e)$.
\end{lemma}
\statusleanthree{Komorebi.PaperIII.EdgeOperators.PacketGraph.EdgeSystem.edgeOp_eq_zero_off_support}

\begin{theorem}[Typed interface annihilation]
\label{thm:edge-annihilation}
If $\tgt(e)\ne\src(f)$, then
\[
  L_fL_e=0.
\]
\end{theorem}
\statusleanthree{Komorebi.PaperIII.EdgeOperators.PacketGraph.EdgeSystem.edgeOp_mul_edgeOp_eq_zero}

\begin{proof}
For a matrix entry $(y,x)$,
\[
  (L_fL_e)_{yx}=\sum_{z\in V}(L_f)_{yz}(L_e)_{zx}.
\]
A potentially nonzero summand would require simultaneously
$z=\src(f)$ and $z=\tgt(e)$, contrary to the hypothesis.
\end{proof}

\begin{theorem}[Matched two-edge collapse]
\label{thm:edge-composable}
If $\tgt(e)=\src(f)$, then
\[
L_fL_e
=
\operatorname{packetMatrix}
\bigl(\src(e),\tgt(f);c(f)c(e)\bigr).
\]
Together with \cref{thm:edge-annihilation}, this gives the complete
composable/noncomposable dichotomy for two edge operators.
\end{theorem}
\statusleanthree{Komorebi.PaperIII.EdgeOperators.PacketGraph.EdgeSystem.edgeOp_mul_edgeOp_of_composable}

\begin{proof}
The matrix-product sum has exactly one potentially nonzero intermediate
index, namely $\tgt(e)=\src(f)$.  At that index the coefficients occur in
operator order as $c(f)c(e)$, and the outer source and target are
$\src(e)$ and $\tgt(f)$.
\end{proof}

\begin{boundary}[Current fiber scope]
The compiled Paper III model above has one scalar coordinate per packet.
For a family of modules $(X_i)_{i\in V}$, an intrinsic edge map would instead
have type $\ell_e:X_{\src(e)}\to X_{\tgt(e)}$, and its totalization on
$\bigoplus_iX_i$ would be
$\iota_{\tgt(e)}\ell_e\pi_{\src(e)}$.  The same support proof applies, but
this dependent fiber-valued generalization is not attributed to the released
scalar Paper III edge-system modules.
\end{boundary}

\Cref{fig:typed-edge-walks} summarizes the support convention and the
composable/noncomposable word dichotomy used throughout the next section.

\begin{figure}[tbp]
\centering
\includegraphics[width=0.98\linewidth]{typed_edge_walks.png}
\caption{The target-row, source-column convention fixes execution order.
Composable words collapse to endpoint blocks, while one broken packet
interface annihilates the entire operator product.}
\label{fig:typed-edge-walks}
\end{figure}

\section{Operator words and directed walks}
\label{sec:words-walks}

\subsection{Typed words}

\begin{definition}[Composable word and execution operator]
\label{def:word-op}
A finite word $w=(e_1,\ldots,e_n)$ is composable when
\[
  \tgt(e_r)=\src(e_{r+1})
  \qquad (1\le r<n).
\]
Set
\[
  L_w=L_{e_n}\cdots L_{e_1},
  \qquad L_{\varnothing}=I.
\]
This definition records the written order of execution even though matrix
composition is read right to left.
\end{definition}
\statusleanthree{Komorebi.PaperIII.EdgeOperators.PacketGraph.EdgeSystem.Composable; Komorebi.PaperIII.EdgeOperators.PacketGraph.EdgeSystem.wordOp}

\begin{theorem}[First broken interface]
\label{thm:first-break}
If a word begins $(e,f,\ldots)$ and $\tgt(e)\ne\src(f)$, then its word
operator is zero.
\end{theorem}
\statusleanthree{Komorebi.PaperIII.EdgeOperators.PacketGraph.EdgeSystem.wordOp_eq_zero_of_first_break}

\begin{theorem}[Nonwalk words vanish]
\label{thm:nonwalk-vanish}
If $w$ is not composable, then $L_w=0$.
\end{theorem}
\statusleanthree{Komorebi.PaperIII.EdgeOperators.PacketGraph.EdgeSystem.wordOp_eq_zero_of_not_composable}

\begin{proof}
Choose the first broken adjacent pair.  The corresponding consecutive factor
is zero by \cref{thm:edge-annihilation}, and associativity absorbs that zero
into the complete product.  The Lean proof implements the same argument by
induction on the list.
\end{proof}

\subsection{The checked recursive expansion}

\begin{definition}[Recursive word expansion]
\label{def:recursive-word-expansion}
Define
\[
  \cW_0=I,
  \qquad
  \cW_{n+1}=T\cW_n.
\]
\end{definition}
\statusleanthree{Komorebi.PaperIII.EdgeOperators.PacketGraph.EdgeSystem.wordExpansion}

\begin{theorem}[Transfer power equals recursive word expansion]
\label{thm:power-recursive-expansion}
For every $n\in\N$,
\[
  T^n=\cW_n.
\]
\end{theorem}
\statusleanthree{Komorebi.PaperIII.EdgeOperators.PacketGraph.EdgeSystem.transferOp_pow_eq_wordExpansion}

\begin{proof}
The zero case is the identity matrix.  If $T^n=\cW_n$, then the
left-successor power convention gives
$T^{n+1}=TT^n=T\cW_n=\cW_{n+1}$.
\end{proof}

\begin{remark}
The formal module also defines an explicit sum over tuples
$w:\operatorname{Fin}(n)\to E$ and a second sum in which noncomposable tuples
are replaced by zero.  It proves that the recursive, explicit, and
walk-filtered matrices agree.  These are separate declarations so that the
finite enumeration claim is auditable rather than hidden in the name
\emph{word expansion}.
\end{remark}
\statusleanthree{Komorebi.PaperIII.EdgeOperators.PacketGraph.EdgeSystem.explicitWordExpansion; Komorebi.PaperIII.EdgeOperators.PacketGraph.EdgeSystem.walkExpansion; Komorebi.PaperIII.EdgeOperators.PacketGraph.EdgeSystem.wordExpansion_eq_explicitWordExpansion; Komorebi.PaperIII.EdgeOperators.PacketGraph.EdgeSystem.explicitWordExpansion_eq_walkExpansion}

\subsection{Exact word-to-walk expansion}

For a nonempty composable word $w=(e_1,\ldots,e_n)$, define its reverse
coefficient product
\[
  c^\leftarrow(w)=c(e_n)c(e_{n-1})\cdots c(e_1).
\]
The reversal is essential over a noncommutative semiring.

\begin{lemma}[One word is one endpoint block]
\label{lem:word-endpoint}
If $w=(e_1,\ldots,e_n)$ is nonempty and composable, then
\[
  L_w
  =
  \operatorname{packetMatrix}
  \bigl(\src(e_1),\tgt(e_n);c^\leftarrow(w)\bigr).
\]
\end{lemma}
\statusleanthree{Komorebi.PaperIII.EdgeOperators.PacketGraph.EdgeSystem.reverseCoefficient; Komorebi.PaperIII.EdgeOperators.PacketGraph.EdgeSystem.terminalTarget; Komorebi.PaperIII.EdgeOperators.PacketGraph.EdgeSystem.wordOp_eq_packetMatrix_of_composable}

\begin{proof}
Induct on $n$.  The case $n=1$ is the definition of $L_e$.  Suppose the
formula holds for the prefix $(e_1,\ldots,e_n)$ and append $f$.  Composability
gives $\tgt(e_n)=\src(f)$.  Multiplying the two packet matrices leaves one
intermediate index and changes the coefficient to
$c(f)c(e_n)\cdots c(e_1)$, with source $\src(e_1)$ and target $\tgt(f)$.
\end{proof}

Write $\Walk_n(i,j)$ for the finite set of length-$n$ composable edge tuples
whose first source is $i$ and whose final target is $j$.  For $n=0$, it has
one empty walk when $i=j$ and is empty otherwise; assign the empty walk the
coefficient $1$.

\begin{theorem}[Finite operator word-to-walk expansion]
\label{thm:word-to-walk}
For every $n\in\N$,
\[
  \boxed{
  T^n
  =
  \sum_{w\in E^n}L_w
  =
  \sum_{w\in E^n}
  \begin{cases}
  L_w,&w\text{ composable},\\
  0,&w\text{ not composable}.
  \end{cases}
  }
\]
Parallel edges occur as distinct summands.
\end{theorem}
\statusleanthree{Komorebi.PaperIII.EdgeOperators.PacketGraph.EdgeSystem.transferOp_pow_eq_explicitWordExpansion; Komorebi.PaperIII.EdgeOperators.PacketGraph.EdgeSystem.transferOp_pow_eq_walkExpansion}

\begin{proof}
Induction and finite distributivity give
\[
\left(\sum_{e\in E}L_e\right)^n
=
\sum_{(e_1,\ldots,e_n)\in E^n}L_{e_n}\cdots L_{e_1}.
\]
Every noncomposable summand is zero by
\cref{thm:nonwalk-vanish}, so inserting the composability filter does not
change the sum.  In Lean, tuple concatenation and the successor sum are
handled by \leanname{tupleWordOp_snoc} and
\leanname{explicitWordExpansion_succ}; the two bridge equalities then yield
the displayed theorem.
\end{proof}

\begin{corollary}[Endpoint coefficient walk sum]
\label{cor:word-to-walk-entry}
For all $n\in\N$ and $i,j\in V$,
\[
  \boxed{
  (T^n)_{ji}
  =
  \sum_{w\in\Walk_n(i,j)}c^\leftarrow(w).
  }
\]
\end{corollary}
\statusleanthree{Komorebi.PaperIII.EdgeOperators.PacketGraph.EdgeSystem.Connects; Komorebi.PaperIII.EdgeOperators.PacketGraph.EdgeSystem.endpointWalkCoefficientSum; Komorebi.PaperIII.EdgeOperators.PacketGraph.EdgeSystem.transferOp_pow_apply_eq_sum_walks}

\begin{proof}
Take the $(j,i)$ entry in \cref{thm:word-to-walk}.  By
\cref{lem:word-endpoint}, a surviving word contributes
$c^\leftarrow(w)$ exactly when its endpoints are $i$ and $j$, and contributes
zero otherwise.  At $n=0$, the identity matrix gives the unique empty
diagonal walk.
\end{proof}

\begin{corollary}[Fixed-length reachability obstruction]
\label{cor:fixed-reachability}
If $\Walk_n(i,j)=\varnothing$, then $(T^n)_{ji}=0$.
\end{corollary}
\statusleanthree{Komorebi.PaperIII.EdgeOperators.PacketGraph.EdgeSystem.transferOp_pow_apply_eq_zero_of_no_walk}

\begin{proof}
This is immediate from the empty sum in the preceding corollary.
\end{proof}

\begin{corollary}[Orbit support obstruction]
\label{cor:orbit-reachability}
Let $S\subseteq V$, and let $x\in R^V$ be supported on $S$.  If $j$ is not
reachable by a directed path from any vertex of $S$, then
\[
  (T^nx)_j=0\qquad\text{for every }n\in\N.
\]
Hence the $T$-orbit of $x$ cannot generate a vector with nonzero
$j$-coordinate.
\end{corollary}
\statusleanthree{Komorebi.PaperIII.EdgeOperators.PacketGraph.EdgeSystem.Reachable; Komorebi.PaperIII.EdgeOperators.PacketGraph.EdgeSystem.transferOp_pow_mulVec_apply_eq_zero_of_unreachable_support}
\statusleanprior{Komorebi.PacketNaturalityCompletion.Reachability.algebraicOrbitHull_le_reachable_support; Komorebi.PacketNaturalityCompletion.Reachability.unreachable_nonzero_packet_obstruction}

\begin{proof}
Expand $(T^nx)_j=\sum_i(T^n)_{ji}x_i$.  Terms with $i\notin S$ vanish by
the support of $x$, and terms with $i\in S$ vanish by
\cref{cor:fixed-reachability}.  The final claim follows because the
$j$-coordinate also vanishes on every finite linear combination of orbit
vectors.
\end{proof}

\section{Canonical exact coordinate transport}
\label{sec:canonical-transport}

\subsection{The Paper III frame model}

Let every chart $a$ carry a linear frame equivalence
$F_a:M\xrightarrow{\sim}M$.  The formal Paper III module defines transport
from $a$ to $b$ by
\[
J_{b\leftarrow a}=F_b^{-1}F_a.
\]
This is the abstract frame form of \emph{synthesize in chart \(a\), then
analyze in chart \(b\)}.

\begin{theorem}[Exact frame transport is coherent]
\label{thm:exact-frame-transport}
For every $a,b,c$,
\[
J_{a\leftarrow a}=\id,\qquad
J_{a\leftarrow b}J_{b\leftarrow a}=\id,\qquad
J_{c\leftarrow b}J_{b\leftarrow a}=J_{c\leftarrow a},
\]
and therefore
\[
J_{a\leftarrow c}J_{c\leftarrow b}J_{b\leftarrow a}=\id.
\]
\end{theorem}
\statusleanthree{Komorebi.PaperIII.Transport.exactTransport_self; Komorebi.PaperIII.Transport.exactTransport_roundTrip; Komorebi.PaperIII.Transport.exactTransport_cocycle; Komorebi.PaperIII.Transport.exactTransport_triangleLoop}

\begin{proof}
Every identity is cancellation of adjacent $F_rF_r^{-1}$ factors.  The Lean
definitions use \leanname{LinearEquiv.trans}; the displayed formulas use the
equivalent right-to-left function-composition notation fixed in this paper.
\end{proof}

\begin{theorem}[Every finite canonical path telescopes]
\label{thm:exact-path-transport}
Let $p=(a_0,a_1,\ldots,a_n)$ be any finite vertex list and put
\[
J_p
=J_{a_n\leftarrow a_{n-1}}\cdots J_{a_1\leftarrow a_0},
\qquad
J_{(a_0)}=\id.
\]
Then
\[
J_p=J_{a_n\leftarrow a_0}.
\]
In particular, if $a_n=a_0$, then $J_p=\id$.
\end{theorem}
\statusleanthree{Komorebi.PaperIII.Transport.pathEndpoint; Komorebi.PaperIII.Transport.exactPathTransport; Komorebi.PaperIII.Transport.exactPathTransport_eq_endpoint; Komorebi.PaperIII.Transport.exactPathTransport_loop}

\begin{proof}
Induct on the list of successive vertices and apply the cocycle in
\cref{thm:exact-frame-transport} at each step.  If the endpoint is the start,
the endpoint transport is the identity.
\end{proof}

\subsection{Exact packet coordinates}

The graph model above resolves a fixed scalar matrix.  Coordinate transport
is a separate construction.  Let $X$ be an $R$-module, and let
\[
\cP^a=(P^a_i)_{i\in I_a}
\]
be a finite complete family of pairwise orthogonal linear idempotents:
\[
P^a_iP^a_j=\delta_{ij}P^a_i,
\qquad
\sum_{i\in I_a}P^a_i=I_X.
\]
Define the exact packet coordinate module
\[
\Coord(\cP^a)
=
\left\{(x_i)_{i\in I_a}:P^a_ix_i=x_i\right\}.
\]
Its analysis and synthesis maps are
\[
\cA_a(x)=(P^a_ix)_i,
\qquad
\cS_a((x_i)_i)=\sum_ix_i.
\]
Completeness and orthogonality give
\[
\cS_a\cA_a=\id_X,
\qquad
\cA_a\cS_a=\id_{\Coord(\cP^a)}.
\]

\begin{definition}[Canonical exact coordinate transport]
\label{def:canonical-J}
For two exact packet systems on the same ambient module, define
\[
  J^X_{b\leftarrow a}
  =
  \cA_b\cS_a:
  \Coord(\cP^a)\longrightarrow\Coord(\cP^b).
\]
The arrow points from the source coordinate system $a$ to the target
coordinate system $b$.
\end{definition}
\statusleanprior{Komorebi.PaperII.NestedRefinement.refineCoordinates; Komorebi.PaperII.NestedRefinement.coordinateRefinementEquiv}

\begin{theorem}[Identity, inverse, and cocycle for canonical transport]
\label{thm:J-cocycle}
For any exact packet systems $a,b,c$ on $X$,
\[
J^X_{a\leftarrow a}=\id,
\qquad
\left(J^X_{b\leftarrow a}\right)^{-1}=J^X_{a\leftarrow b},
\qquad
J^X_{c\leftarrow b}J^X_{b\leftarrow a}=J^X_{c\leftarrow a}.
\]
\end{theorem}
\statusleanprior{Komorebi.PaperII.NestedRefinement.refineCoordinates_self; Komorebi.PaperII.NestedRefinement.coarsenCoordinates_refineCoordinates; Komorebi.PaperII.NestedRefinement.refineCoordinates_coarsenCoordinates; Komorebi.PaperII.NestedRefinement.refineCoordinates_comp; Komorebi.PaperII.NestedRefinement.coordinateRefinementEquiv_trans}

\begin{proof}
Insert $\cS_b\cA_b=\id_X$:
\[
J^X_{c\leftarrow b}J^X_{b\leftarrow a}
=\cA_c\cS_b\cA_b\cS_a
=\cA_c\cS_a
=J^X_{c\leftarrow a}.
\]
The identity statement is $\cA_a\cS_a=\id$.  Setting $c=a$ in the cocycle
gives $J^X_{a\leftarrow b}J^X_{b\leftarrow a}=\id$; interchanging $a$ and $b$
gives the other composition, so the two transports are inverses.
\end{proof}

If $F:X\to Y$ is linear and $\cP^a,\cP^b$ are source packet systems while
$\cQ^r,\cQ^s$ are target packet systems, put
\[
\widehat F_{r\leftarrow a}
=
\cA^Y_rF\cS^X_a.
\]

\begin{theorem}[Canonical coordinate naturality]
\label{thm:canonical-naturality}
The exact coordinate representation preserves typed composition and satisfies
\[
J^Y_{s\leftarrow r}\widehat F_{r\leftarrow a}
=
\widehat F_{s\leftarrow b}J^X_{b\leftarrow a}.
\]
\end{theorem}
\statusleanprior{Komorebi.PaperII.NestedRefinement.coordinateOperator_comp; Komorebi.PaperII.NestedRefinement.coordinateOperator_naturality}

\begin{proof}
Both sides reduce to $\cA^Y_sF\cS^X_a$ after inserting the appropriate
synthesis--analysis identity.  The composition law is the same calculation
with the middle ambient module inserted.
\end{proof}

\begin{corollary}[Canonical triangle holonomy is trivial]
\label{cor:J-flat}
For every triple $a,b,c$,
\[
J^X_{a\leftarrow c}J^X_{c\leftarrow b}J^X_{b\leftarrow a}=\id.
\]
\end{corollary}
\statusleanprior{Komorebi.PaperII.NestedRefinement.refineCoordinates_comp; Komorebi.PaperII.NestedRefinement.refineCoordinates_self}

\begin{boundary}[Refinement is not reduction]
Although Paper II uses the name \emph{refineCoordinates}, $J$ is an
equivalence between two complete coordinate presentations.  It neither
truncates packets nor projects to a smaller state space.  Any lossy
restriction, interpolation, or learned identification is extra data and
cannot inherit \cref{thm:J-cocycle} without proof.
\end{boundary}

\section{A distinct noncanonical connection API}
\label{sec:connection}

\subsection{The minimal compiled API}

The formal Paper III API starts with deliberately weak assumptions.  All
chart transports act on one module $M$, so chart labels distinguish routes
but not Lean types.

\begin{definition}[Self-normalized noncanonical connection]
\label{def:minimal-connection}
A minimal connection is a family
\[
\Gamma_{b\leftarrow a}\in\operatorname{Aut}_R(M)
\qquad(a,b\in B)
\]
with only
\[
\Gamma_{a\leftarrow a}=\id_M.
\]
Neither reverse-edge inversion nor the cocycle law is a field of the
structure.
\end{definition}
\statusleanthree{Komorebi.PaperIII.Transport.Connection}

At a triple $(a,b,c)$, the compiled cocycle predicate is
\[
\Gamma_{c\leftarrow b}\Gamma_{b\leftarrow a}
=\Gamma_{c\leftarrow a}.
\]
Flatness means this equality holds for every triple.  The formal
multiplicative relative defect and triangle holonomy are
\[
D^\times_{c,b,a}
=
\bigl(\Gamma_{c\leftarrow b}\Gamma_{b\leftarrow a}\bigr)
\Gamma_{c\leftarrow a}^{-1},
\qquad
H_{c,b,a}
=
\Gamma_{a\leftarrow c}
\Gamma_{c\leftarrow b}
\Gamma_{b\leftarrow a}.
\]
The defect is based at the target $c$; the holonomy is based at $a$.
\statusleanthree{Komorebi.PaperIII.Transport.Connection.CocycleAt; Komorebi.PaperIII.Transport.Connection.Flat; Komorebi.PaperIII.Transport.Connection.cocycleDefect; Komorebi.PaperIII.Transport.Connection.triangleHolonomy}

\begin{theorem}[Checked flatness consequences]
\label{thm:minimal-flatness}
For the minimal connection API, flatness implies
\[
D^\times_{c,b,a}=\id_M,
\qquad
H_{c,b,a}=\id_M
\]
for every $a,b,c$.
\end{theorem}
\statusleanthree{Komorebi.PaperIII.Transport.Connection.cocycleDefect_eq_refl_of_flat; Komorebi.PaperIII.Transport.Connection.triangleHolonomy_eq_refl_of_flat}

\begin{proof}
Flatness replaces the two-edge route in $D^\times$ by the direct route,
which cancels its inverse.  For $H$, flatness first collapses $a\to b\to c$
to $a\to c$, then collapses $a\to c\to a$ to the normalized transport
$a\to a$.
\end{proof}

\begin{proposition}[The minimal API admits nontrivial holonomy]
\label{prop:twisted-loop}
There is a connection on three points over $\mathbb Q$ for which transport
$0\to1$ is multiplication by $2$, every other unequal-endpoint transport is
the identity, and
\[
H_{2,1,0}(1)=2.
\]
In particular $H_{2,1,0}\ne\id$.
\end{proposition}
\statusleanthree{Komorebi.PaperIII.Transport.twistedThreePointConnection; Komorebi.PaperIII.Transport.twisted_triangleHolonomy_apply_one; Komorebi.PaperIII.Transport.twisted_triangleHolonomy_ne_refl}

\begin{boundary}[Why the minimal defect and loop are distinct]
Without reverse-edge inversion,
$\Gamma_{a\leftarrow c}$ need not equal
$\Gamma_{c\leftarrow a}^{-1}$.  Therefore the multiplicative relative defect
$D^\times_{c,b,a}$ and triangle holonomy $H_{c,b,a}$ are different
comparisons.  The formal module proves that global flatness trivializes both;
it does not identify them for an arbitrary connection.
\end{boundary}

\begin{figure}[H]
\centering
\includegraphics[width=0.98\linewidth]{transport_vs_connection.png}
\caption{Exact coordinate transport and connection transport are distinct.
The canonical loop telescopes to the identity.  A chosen connection may
carry nontrivial holonomy, which transforms by conjugation under gauge.}
\label{fig:transport-connection}
\end{figure}

\Cref{fig:transport-connection} records this separation visually: the exact
coordinate route telescopes, whereas the chosen connection route may carry a
nonidentity loop.

\subsection{A stronger reversible pair API}

Let $R$ be a commutative ring, let $B$ be a type of charts or refinement
states, and let each $C_a$ be an $R$-module.  The chart type is not assumed
finite.  The following definition is deliberately not derived from analysis
and synthesis.

\begin{definition}[Normalized reversible connection]
\label{def:connection}
A normalized reversible connection is a family of linear equivalences
\[
\Gamma_{b\leftarrow a}:C_a\xrightarrow{\;\sim\;}C_b
\qquad(a,b\in B)
\]
such that
\[
\Gamma_{a\leftarrow a}=\id,
\qquad
\Gamma_{a\leftarrow b}=\Gamma_{b\leftarrow a}^{-1}.
\]
No cocycle equation is included in the structure.
\end{definition}
\statusleanthree{Komorebi.PaperIII.Transport.ReversibleConnection}

This strengthens the homogeneous minimal connection with heterogeneous
fibers and reverse-edge inversion.  It remains additional data, not an alias
for canonical coordinate transport.

\begin{definition}[Calibration and flatness]
\label{def:calibration-flatness}
When the $C_a$ are exact packet coordinates of one ambient module, the
connection is \emph{canonically calibrated} if
\[
\Gamma_{b\leftarrow a}=J_{b\leftarrow a}
\quad\text{for every }a,b.
\]
It is \emph{flat} if
\[
\Gamma_{c\leftarrow b}\Gamma_{b\leftarrow a}
=
\Gamma_{c\leftarrow a}
\quad\text{for every }a,b,c.
\]
The canonical connection is flat; a general flat connection is proved in
\cref{thm:flat-gauge-canonical} to be gauge-equivalent to it.
\end{definition}
\statusleanthree{Komorebi.PaperIII.Transport.ReversibleConnection.CocycleAt; Komorebi.PaperIII.Transport.ReversibleConnection.Flat; Komorebi.PaperIII.Transport.ReversibleConnection.canonicalConnection; Komorebi.PaperIII.Transport.ReversibleConnection.canonicalConnection_flat}

The words \emph{calibrated} and \emph{flat} have only the algebraic meanings
displayed above.

\begin{remark}[Graph-local variant]
One may specify $\Gamma_e$ only on the oriented edges of a chart graph,
together with inverse maps on reversed edges.  Path transport and loop
holonomy below still make sense.  A direct cocycle defect for $(a,b,c)$,
however, requires a declared direct transport from $a$ to $c$.  The pairwise
API in \cref{def:connection} keeps this comparison total.
\end{remark}

\subsection{Cocycle defect and triangle holonomy}

The commutative-ring hypothesis makes subtraction of linear maps available
from this point through \cref{sec:obstruction-packets}.

\begin{definition}[Cocycle defect]
\label{def:cocycle-defect}
For $a,b,c\in B$, define the additive cocycle defect
\[
\Delta_{c,b,a}
=
\Gamma_{c\leftarrow b}\Gamma_{b\leftarrow a}
-\Gamma_{c\leftarrow a}
\in\Hom_R(C_a,C_c).
\]
\end{definition}
\statusleanthree{Komorebi.PaperIII.Transport.ReversibleConnection.cocycleDefect}

The manuscript subscript order is target, intermediate, source.  The Lean
arguments are supplied in traversal order \leanname{(a,b,c)}.

\begin{definition}[Triangle holonomy]
\label{def:triangle-holonomy}
The triangle holonomy based at $a$ is
\[
H_{c,b,a}
=
\Gamma_{a\leftarrow c}
\Gamma_{c\leftarrow b}
\Gamma_{b\leftarrow a}
\in\operatorname{Aut}_R(C_a).
\]
It traverses $a\to b\to c\to a$.
\end{definition}
\statusleanthree{Komorebi.PaperIII.Transport.ReversibleConnection.triangleHolonomy; Komorebi.PaperIII.Transport.ReversibleConnection.triangleObstruction}

Holonomy is attached to the chosen connection $\Gamma$, never to an
unspecified coordinate change.

\begin{theorem}[Defect--holonomy identity]
\label{thm:defect-holonomy}
For every triple $(a,b,c)$,
\[
\boxed{
H_{c,b,a}-\id_{C_a}
=
\Gamma_{a\leftarrow c}\Delta_{c,b,a}.
}
\]
Consequently,
\[
\Delta_{c,b,a}=0
\quad\Longleftrightarrow\quad
H_{c,b,a}=\id_{C_a}.
\]
\end{theorem}
\statusleanthree{Komorebi.PaperIII.Transport.ReversibleConnection.triangleHolonomy_sub_id; Komorebi.PaperIII.Transport.ReversibleConnection.cocycleDefect_eq_zero_iff_triangleHolonomy_eq_refl}

\begin{proof}
Using reversibility,
\[
\begin{aligned}
\Gamma_{a\leftarrow c}\Delta_{c,b,a}
&=
\Gamma_{a\leftarrow c}
\Gamma_{c\leftarrow b}
\Gamma_{b\leftarrow a}
-
\Gamma_{a\leftarrow c}\Gamma_{c\leftarrow a}\\
&=H_{c,b,a}-\id_{C_a}.
\end{aligned}
\]
Left composition by the equivalence $\Gamma_{a\leftarrow c}$ is injective,
which gives the equivalence of the two vanishing statements.
\end{proof}

\begin{corollary}[Canonical transport has zero connection defect]
\label{cor:canonical-zero-defect}
If $\Gamma_{b\leftarrow a}=J_{b\leftarrow a}$, then
\[
\Delta_{c,b,a}=0,
\qquad
H_{c,b,a}=\id
\]
for all triples.
\end{corollary}
\statusleanthree{Komorebi.PaperIII.Transport.ReversibleConnection.canonicalTransport_zero_connectionDefect; Komorebi.PaperIII.Transport.ReversibleConnection.canonicalTransport_triangleHolonomy_eq_refl}

\begin{proof}
The defect vanishes by the canonical cocycle in
\cref{thm:J-cocycle}; apply \cref{thm:defect-holonomy}.
\end{proof}

\subsection{Paths and endpoint independence}

\begin{definition}[Path-ordered connection transport]
\label{def:path-transport}
For a vertex sequence
$p=(a_0,a_1,\ldots,a_n)$, define
\[
\Gamma_p
=
\Gamma_{a_n\leftarrow a_{n-1}}\cdots
\Gamma_{a_1\leftarrow a_0},
\qquad
\Gamma_{(a_0)}=\id_{C_{a_0}}.
\]
If $a_n=a_0$, then $\Gamma_p$ is the holonomy of the loop $p$.
\end{definition}
\statusleanthree{Komorebi.PaperIII.Transport.ReversibleConnection.pathTransport; Komorebi.PaperIII.Transport.ReversibleConnection.VertexPath; Komorebi.PaperIII.Transport.ReversibleConnection.indexedPathTransport}

Products are path-ordered: the first traversed edge acts on the right.

\begin{proposition}[Flat pair transport collapses to endpoints]
\label{prop:flat-endpoints}
If $\Gamma$ is flat, then every path satisfies
\[
\Gamma_p=\Gamma_{a_n\leftarrow a_0}.
\]
In particular, every loop has identity holonomy.
\end{proposition}
\statusleanthree{Komorebi.PaperIII.Transport.ReversibleConnection.pathTransport_eq_direct_of_flat; Komorebi.PaperIII.Transport.ReversibleConnection.indexedPathTransport_eq_direct_of_flat; Komorebi.PaperIII.Transport.ReversibleConnection.indexedLoopTransport_eq_refl_of_flat}

\begin{proof}
Induct on the path length.  Appending $a_{n+1}$ left-multiplies the
inductive expression by $\Gamma_{a_{n+1}\leftarrow a_n}$, and flatness
replaces the two factors by
$\Gamma_{a_{n+1}\leftarrow a_0}$.  For a loop, normalization gives
$\Gamma_{a_0\leftarrow a_0}=\id$.
\end{proof}

\begin{remark}
The converse \emph{trivial holonomy implies flatness} depends on which loops
are present in a graph-local API.  In the complete pair API, triviality of
every triangle holonomy is equivalent to flatness by
\cref{thm:defect-holonomy}.  No topological simple-connectivity theorem is
being invoked.
\end{remark}

\section{Gauge covariance and comparison with canonical transport}
\label{sec:gauge}

\begin{definition}[Gauge change]
\label{def:gauge-change}
A gauge is a family of automorphisms $G_a\in\operatorname{Aut}_R(C_a)$.
It acts on a connection by
\[
\Gamma^G_{b\leftarrow a}
=
G_b\Gamma_{b\leftarrow a}G_a^{-1}.
\]
\end{definition}
\statusleanthree{Komorebi.PaperIII.Transport.ReversibleConnection.Gauge; Komorebi.PaperIII.Transport.ReversibleConnection.gauge}

A gauge changes coordinate representatives at each chart; it does not change
the chart graph.

\begin{theorem}[Gauge covariance]
\label{thm:gauge-covariance}
The defect and triangle holonomy transform as
\[
\Delta^G_{c,b,a}
=G_c\Delta_{c,b,a}G_a^{-1},
\qquad
H^G_{c,b,a}
=G_aH_{c,b,a}G_a^{-1}.
\]
Thus flatness and triviality of triangle holonomy are gauge invariant.
\end{theorem}
\statusleanthree{Komorebi.PaperIII.Transport.ReversibleConnection.gauge_cocycleDefect; Komorebi.PaperIII.Transport.ReversibleConnection.gauge_triangleHolonomy; Komorebi.PaperIII.Transport.ReversibleConnection.gauge_flat_iff; Komorebi.PaperIII.Transport.ReversibleConnection.gauge_triangleHolonomy_eq_refl_iff}

\begin{proof}
Substitute the definition of $\Gamma^G$.  Adjacent factors
$G_b^{-1}G_b$ and $G_c^{-1}G_c$ cancel.  In the defect, both the composite
and direct terms acquire the same outer factors $G_c$ and $G_a^{-1}$.  In
the loop, all inner factors cancel and only conjugation at the base remains.
\end{proof}

\begin{theorem}[Flat connections are gauges of the canonical pair transport]
\label{thm:flat-gauge-canonical}
Suppose the $C_a$ are exact packet-coordinate modules for one ambient module,
so the canonical $J_{b\leftarrow a}$ is available.  If $\Gamma$ is a flat
normalized reversible pair connection, choose $o\in B$ and set
\[
G_a=\Gamma_{a\leftarrow o}J_{o\leftarrow a}.
\]
Then every pair transport satisfies
\[
\boxed{
\Gamma_{b\leftarrow a}
=G_bJ_{b\leftarrow a}G_a^{-1}.
}
\]
\end{theorem}
\statusleanthree{Komorebi.PaperIII.Transport.ReversibleConnection.flatGauge; Komorebi.PaperIII.Transport.ReversibleConnection.flat_isGaugeOfCanonical}

\begin{proof}
Reversibility gives
$G_a^{-1}=J_{a\leftarrow o}\Gamma_{o\leftarrow a}$.  Therefore
\[
\begin{aligned}
G_bJ_{b\leftarrow a}G_a^{-1}
&=
\Gamma_{b\leftarrow o}
J_{o\leftarrow b}J_{b\leftarrow a}
J_{a\leftarrow o}
\Gamma_{o\leftarrow a}\\
&=
\Gamma_{b\leftarrow o}\Gamma_{o\leftarrow a}\\
&=
\Gamma_{b\leftarrow a}.
\end{aligned}
\]
The first reduction uses the canonical cocycle twice; the second uses
flatness of $\Gamma$.
\end{proof}

\begin{boundary}[Flat does not mean calibrated]
The theorem produces a gauge relating $\Gamma$ to $J$.  It does not prove
$\Gamma=J$ before that gauge is chosen.  Conversely, nontrivial triangle
holonomy cannot be removed by a gauge, because conjugation cannot turn a
nonidentity automorphism into the identity.
\end{boundary}

\section{Obstruction packets: local certificates and holonomy blocks}
\label{sec:obstruction-packets}

\subsection{The generic compiled certificate}

Let $\mathsf{Candidate}$, $\mathsf{Packet}$, and $\mathsf{Residual}$ be
arbitrary types, with a distinguished
$0\in\mathsf{Residual}$, and let
\[
r:\mathsf{Candidate}\longrightarrow
\mathsf{Packet}\longrightarrow\mathsf{Residual}
\]
be a residual function.

\begin{definition}[Candidate-specific solution and obstruction packet]
\label{def:generic-obstruction}
A candidate $x$ is a solution when
\[
\operatorname{IsSolution}_r(x)
\quad\Longleftrightarrow\quad
\forall p,\ r(x,p)=0.
\]
An obstruction packet for $x$ is a pair consisting of a named packet $p$
and a certificate $r(x,p)\ne0$.
\end{definition}
\statusleanthree{Komorebi.PaperIII.NonlinearObstructions.IsSolution; Komorebi.PaperIII.NonlinearObstructions.ObstructionPacket}

\begin{theorem}[Exact logical force of a local obstruction]
\label{thm:generic-obstruction}
Every obstruction packet for $x$ proves that $x$ is not a solution, and
\[
\neg\operatorname{IsSolution}_r(x)
\quad\Longleftrightarrow\quad
\text{there exists an obstruction packet for }x.
\]
The reverse implication from failure to a packet witness is classical and
does not require the packet type to be finite.
\end{theorem}
\statusleanthree{Komorebi.PaperIII.NonlinearObstructions.ObstructionPacket.not_isSolution; Komorebi.PaperIII.NonlinearObstructions.obstructionPacket_of_not_isSolution; Komorebi.PaperIII.NonlinearObstructions.not_isSolution_iff_nonempty_obstructionPacket}

\begin{proof}
A purported solution would make the residual vanish at the certified packet,
contradicting its nonzero field.  Conversely, if no packet carries a nonzero
residual, then every residual is zero and the candidate satisfies the
definition of solution.
\end{proof}

\begin{boundary}[Candidate-specific means candidate-specific]
The theorem rules out the named $x$.  It neither constructs a solution nor
shows that every candidate is obstructed.  It provides no minimality,
stability, computability, or classification theorem for the packet witness.
\end{boundary}

\subsection{Holonomy-block specialization}

For a reversible connection, the generic residual may be specialized at a
fixed triangle by taking a packet label to be a pair $(\mu,\nu)$ and the
residual to be the corresponding block of $H-\id$.  A block is therefore a
packet residual; when it is nonzero, the label together with a proof of
nonvanishing is an obstruction-packet certificate in the sense of
\cref{def:generic-obstruction}.  The reconstruction theorem below is the extra
statement needed to turn vanishing of all those local residuals into the
cocycle equation at that triangle.

Fix a triangle $(a,b,c)$ and resolve the base fiber $C_a$ by a finite complete
family of idempotents $(\Pi^a_\mu)_{\mu\in I_a}$.  Orthogonality may also hold,
as it does for the exact range-coordinate projectors below, but the compiled
reconstruction theorem does not require it.  In exact range coordinates,
$\Pi^a_\mu$ is simply the projector that retains component $\mu$ and zeros
the other components.

\begin{definition}[Holonomy obstruction and packet residuals]
\label{def:obstruction-packet}
Define
\[
\Ob_{c,b,a}=H_{c,b,a}-\id_{C_a},
\qquad
\Ob^{\mu\nu}_{c,b,a}
=
\Pi^a_\mu\Ob_{c,b,a}\Pi^a_\nu.
\]
The map $\Ob^{\mu\nu}_{c,b,a}$ is the obstruction block, or packet residual,
from input packet $\nu$ to output packet $\mu$.  If it is nonzero, then
$(\mu,\nu)$ and that nonvanishing proof form an obstruction-packet
certificate.  Diagonal blocks detect within-packet holonomy; off-diagonal
blocks detect holonomy-induced packet mixing.
\end{definition}
\statusleanthree{Komorebi.PaperIII.HolonomyBlocks.PacketResolution; Komorebi.PaperIII.HolonomyBlocks.PacketResolution.block}

An obstruction-packet certificate here is a nonzero block of a precisely
specified holonomy defect together with its packet-pair label; it is not
merely a heuristic label for a large residual.

\begin{theorem}[Exact obstruction-block reconstruction and detection]
\label{thm:obstruction-detection}
The total obstruction has the exact finite resolution
\[
\boxed{
\Ob_{c,b,a}
=
\sum_{\mu,\nu\in I_a}
\Ob^{\mu\nu}_{c,b,a}.
}
\]
Moreover, the following are equivalent:
\begin{enumerate}[label=(\roman*)]
\item $\Delta_{c,b,a}=0$;
\item $H_{c,b,a}=\id$;
\item $\Ob_{c,b,a}=0$;
\item $\Ob^{\mu\nu}_{c,b,a}=0$ for every $\mu,\nu$.
\end{enumerate}
\end{theorem}
\statusleanthree{Komorebi.PaperIII.HolonomyBlocks.obstruction_sum_blocks; Komorebi.PaperIII.HolonomyBlocks.obstructionBlocks_zero_iff; Komorebi.PaperIII.HolonomyBlocks.obstructionDetection_equivalences}

\begin{proof}
Insert completeness on both sides:
\[
\Ob
=
\left(\sum_\mu\Pi^a_\mu\right)
\Ob
\left(\sum_\nu\Pi^a_\nu\right)
=
\sum_{\mu,\nu}\Pi^a_\mu\Ob\Pi^a_\nu.
\]
Thus (iv) implies (iii), and (iii) implies (iv) by composition with the
projectors.  Conditions (ii) and (iii) are equivalent by the definition of
$\Ob$ and additive cancellation.
Conditions (i) and (ii) are equivalent by
\cref{thm:defect-holonomy}.
\end{proof}

\Cref{fig:holonomy-obstruction-blocks} depicts the two-sided finite block
resolution behind the reconstruction identity and its off-diagonal witnesses.

\begin{figure}[H]
\centering
\includegraphics[width=0.98\linewidth]{holonomy_obstruction_blocks.png}
\caption{A finite complete packet resolution reconstructs the total holonomy
obstruction from its input-output blocks.  A nonzero off-diagonal block is an
exact packet-mixing witness.}
\label{fig:holonomy-obstruction-blocks}
\end{figure}

\begin{corollary}[Packet-supported witness]
\label{cor:packet-witness}
If $\Ob^{\mu\nu}_{c,b,a}\ne0$, then there is a vector $x$ supported in packet
$\nu$ such that
\[
\Pi^a_\mu\Ob_{c,b,a}x\ne0.
\]
Such an $x$ is an exact finite witness of nonflatness.
\end{corollary}
\statusleanthree{Komorebi.PaperIII.HolonomyBlocks.exists_packetWitness_of_obstructionBlock_ne_zero}

\begin{proof}
Classically choose $y$ on which the nonzero linear map
$\Pi^a_\mu\Ob\Pi^a_\nu$ is nonzero and set $x=\Pi^a_\nu y$.
Idempotence gives $\Pi^a_\nu x=x$, and the required output is the chosen
nonzero value.
\end{proof}

\subsection{Gauge behavior of packet blocks}

Vanishing or nonvanishing of the total triangle obstruction is gauge invariant
by \cref{thm:gauge-covariance}, but a fixed packet pattern need not be.  A
general $G_a$ can mix the projectors $\Pi^a_\mu$.

\begin{proposition}[Packet-preserving gauge covariance]
\label{prop:packet-gauge}
If $G_a\Pi^a_\mu=\Pi^a_\mu G_a$ for every $\mu$, then
\[
\left(\Ob^G\right)^{\mu\nu}
=
G_a\Ob^{\mu\nu}G_a^{-1}.
\]
Consequently the zero/nonzero status of each obstruction block is preserved
by packet-preserving gauges.
\end{proposition}
\statusleanthree{Komorebi.PaperIII.HolonomyBlocks.obstructionBlock_gauge_of_commutes; Komorebi.PaperIII.HolonomyBlocks.obstructionBlock_gauge_eq_zero_iff}

\begin{proof}
Use $\Ob^G=G_a\Ob G_a^{-1}$ and commute $G_a$ and $G_a^{-1}$ through the
two packet projectors.  Composition by an equivalence is injective in both
directions.
\end{proof}

\begin{boundary}[Packet patterns are not absolute]
Without the commutation hypothesis in \cref{prop:packet-gauge}, only the
vanishing of the total obstruction is gauge invariant.  The location of
individual nonzero blocks may change under a gauge that mixes packets.
\end{boundary}

\section{Packet-supported products and holonomy walks}
\label{sec:holonomy-walks}

The word-to-walk mechanism is not special to powers of one matrix.  Let $V$
be one finite packet type, let $R$ be a semiring, and let
$B_1,\ldots,B_m\in R^{V\times V}$ be a finite sequence of possibly different
packet matrices.  Rows are targets and columns are sources as before.

\begin{proposition}[Heterogeneous packet-path expansion]
\label{prop:heterogeneous-path}
For $i,j\in V$ and $m\geq1$, set $i_0=i$ and $i_m=j$.  Then
\[
\boxed{
(B_m\cdots B_1)_{ji}
=
\sum_{(i_1,\ldots,i_{m-1})\in V^{m-1}}
  B_m[i_m,i_{m-1}]\cdots B_1[i_1,i_0].
}
\]
When $m=1$, the sum is over the unique empty tuple and its summand is
$B_1[j,i]$.
For $m=0$, the product is the identity and the entry is the Kronecker delta
$\delta_{ji}$.  If every $B_r[q,p]$ vanishes outside a declared relation
$\mathcal R_r(p,q)$, then the endpoint entry is zero whenever there is no
packet path compatible with the successive relations.
\end{proposition}
\statusleanthree{Komorebi.PaperIII.EdgeOperators.heterogeneousProduct; Komorebi.PaperIII.EdgeOperators.heterogeneousPathSum; Komorebi.PaperIII.EdgeOperators.heterogeneousPacketPathExpansion; Komorebi.PaperIII.EdgeOperators.SupportedStep; Komorebi.PaperIII.EdgeOperators.SupportedPath; Komorebi.PaperIII.EdgeOperators.heterogeneousProduct_apply_eq_zero_of_no_supportedPath}

\begin{proof}
Expand each matrix product entry and distribute the finite sums.  Induction on
the factor list yields the displayed path sum, with the empty list evaluated
as the identity matrix.  If no supported intermediate packet exists, every
summand contains an entry that is zero by its declared support relation.
\end{proof}

\begin{boundary}[Changing packet types along a chain]
The compiled theorems
\leanname{heterogeneousPacketPathExpansion_two} and
\leanname{heterogeneousPacketPathExpansion_three} allow different packet
types at the successive stages of two- and three-factor products.  The
arbitrary-length theorem above permits different factors but uses one common
finite packet type $V$.  A fully dependent arbitrary-length chain of changing
modules and packet types remains outside the version 1.0.1 contract.
\end{boundary}

\begin{remark}[Checked specialization]
When every $B_r$ is one repeated endomorphism $T$, the recursive packet-power
specialization is inherited from Paper I.
\end{remark}
\statusleanprior{Komorebi.PacketNaturalityCompletion.Packets.transferBlock_pow_eq_packetPathSum}

For a loop represented on one common finite scalar packet-coordinate type,
the proposition resolves its successive transport matrices into packet
walks.  If $R$ is a ring, as in the reversible-connection API, then
$(H-I)_{ji}=H_{ji}-\delta_{ji}$; in particular, off-diagonal entries are
unchanged.  Thus a nonzero off-diagonal holonomy obstruction in this
common-index model requires a transport word along the closed chart loop
whose packet path runs from its input packet to its output packet.  This
conclusion is conditional on declared block supports for the connection
matrices; an arbitrary connection need not respect the transfer graph from
\cref{sec:edge-operators}.

\section{A minimal exact example}
\label{sec:example}

\begin{example}[A scalar triangle]
Let $B=\{0,1,2\}$ and $C_0=C_1=C_2=\bbk$ for a field $\bbk$.  Choose
$u,v,w\in\bbk^\times$ and set
\[
\Gamma_{1\leftarrow0}(x)=ux,\qquad
\Gamma_{2\leftarrow1}(x)=vx,\qquad
\Gamma_{2\leftarrow0}(x)=wx,
\]
with reverse transports given by inverse scalars.  Then
\[
\Delta_{2,1,0}(x)=(vu-w)x,
\]
and the triangle holonomy based at $0$ is
\[
H_{2,1,0}(x)=w^{-1}vux.
\]
Therefore
\[
H_{2,1,0}-\id
=w^{-1}\Delta_{2,1,0}.
\]
The connection is flat on this triangle exactly when $vu=w$.  If the three
coordinate systems are literally the same and canonical transport is used,
then $u=v=w=1$ and the obstruction vanishes.  Choosing arbitrary nonzero
$u,v,w$ is a connection choice, not a consequence of changing exact packet
coordinates.
\end{example}

\begin{example}[A two-step transfer walk]
Let a packet graph have edges $e:0\to1$ and $f:1\to2$, with no other edges.
Then
\[
T=L_e+L_f,\qquad
(T^2)_{20}=c(f)c(e).
\]
The products $L_eL_e$, $L_eL_f$, and $L_fL_f$ vanish by typed mismatch;
only $L_fL_e$, corresponding to the walk $0\to1\to2$, contributes to the
$(2,0)$ entry.  This example also fixes the noncommutative coefficient order.
\end{example}

\section{Canonical coordinates for arbitrary maps}
\label{sec:arbitrary-maps}

The analysis and synthesis identities are identities of functions, not only
of linear maps.  Consequently exact coordinate change does not become
curved merely because the represented operation is nonlinear.

\begin{proposition}[Set-level exact naturality]
\label{prop:set-naturality}
Let $\cA^X_a:X\xrightarrow{\sim}X_a$ and
$\cA^X_b:X\xrightarrow{\sim}X_b$ be exact coordinate equivalences for $X$,
with inverses $\cS^X_a$ and $\cS^X_b$.  Likewise, let
$\cA^Y_r:Y\xrightarrow{\sim}Y_r$ and
$\cA^Y_s:Y\xrightarrow{\sim}Y_s$ be exact coordinate equivalences for $Y$,
with inverses $\cS^Y_r$ and $\cS^Y_s$.  Let $F:X\to Y$ be any function, with
no linearity or differentiability hypothesis.  Set
\[
J^X_{b\leftarrow a}=\cA^X_b\circ\cS^X_a,
\qquad
J^Y_{s\leftarrow r}=\cA^Y_s\circ\cS^Y_r,
\]
and define
\[
\widehat F_{r\leftarrow a}
=
\cA^Y_r\circ F\circ\cS^X_a.
\]
Then
\[
J^Y_{s\leftarrow r}\circ\widehat F_{r\leftarrow a}
=
\widehat F_{s\leftarrow b}\circ J^X_{b\leftarrow a}.
\]
If $\cA^Z_t:Z\xrightarrow{\sim}Z_t$ is another exact coordinate equivalence
and $G:Y\to Z$ is any function, then exact coordinate representation also
preserves composition:
\[
\widehat{G\circ F}_{t\leftarrow a}
=
\widehat G_{t\leftarrow r}\circ\widehat F_{r\leftarrow a}.
\]
\end{proposition}
\statusleanthree{Komorebi.PaperIII.Transport.exactCoordinateMap; Komorebi.PaperIII.Transport.exactCoordinateMap_comp; Komorebi.PaperIII.Transport.exactCoordinateMap_naturality}

\begin{proof}
Both sides of the naturality equation reduce to
$\cA^Y_s\circ F\circ\cS^X_a$ using
$\cS\circ\cA=\id$ in $X$ and $Y$.  For composition, the middle
$\cS^Y\circ\cA^Y$ cancels in the same way.
\end{proof}

\begin{boundary}[Where nonlinear defects can enter]
A nonlinear \emph{coordinate} defect cannot be blamed on the exact canonical
$J$ of \cref{def:canonical-J}; \cref{prop:set-naturality} makes that defect
identically zero.  Nonzero defects require extra structure such as a
noncanonical connection, truncation, approximation, restricted domains, or
failure of an operation to respect a declared packet model.
\end{boundary}

\section{Nonlinear operation trees: proved locality and scoped roadmap}
\label{sec:operation-trees}

\subsection{The compiled unary--binary foundation}

\begin{definition}[Finite operation-tree grammar]
\label{def:compiled-operation-tree}
For types of input labels $I$, constants $K$, unary symbols $U$, and binary
symbols $B$, the released formal grammar has four constructors:
\[
\begin{aligned}
t ::={}&\operatorname{input}(i)
\mid\operatorname{constant}(k)
\mid\operatorname{unary}(u,t)\\
&\mid\operatorname{binary}(b,t_1,t_2).
\end{aligned}
\]
Given an arbitrary value type $Y$, evaluation takes completely explicit maps
\[
I\to Y,\qquad K\to Y,\qquad U\to Y\to Y,\qquad
B\to Y\to Y\to Y.
\]
No algebraic, continuity, or differentiability law is imposed on these
semantics.
\end{definition}
\statusleanthree{Komorebi.PaperIII.NonlinearObstructions.OperationTree; Komorebi.PaperIII.NonlinearObstructions.OperationTree.eval}

When $I$ has decidable equality, let $\operatorname{inputs}(t)$ be the finite
set of input labels actually occurring in $t$.

\begin{theorem}[Operation-tree input locality]
\label{thm:tree-locality}
Let $\eta_1,\eta_2:I\to Y$ be two input environments.  If
\[
\eta_1(i)=\eta_2(i)
\qquad\text{for every }i\in\operatorname{inputs}(t),
\]
then the two evaluations of $t$ are equal under the same constant, unary, and
binary semantics.  In particular, if
$\operatorname{inputs}(t)=\varnothing$, evaluation is independent of the
input environment.
\end{theorem}
\statusleanthree{Komorebi.PaperIII.NonlinearObstructions.OperationTree.eval_congr_on_inputs; Komorebi.PaperIII.NonlinearObstructions.OperationTree.eval_eq_of_inputs_eq_empty}

\begin{proof}
Structural induction on $t$.  The input case is the agreement hypothesis;
the constant case is reflexive; the unary case applies the induction
hypothesis to its child; and the binary case applies it independently to the
left and right subtrees, whose input sets form a union.
\end{proof}

\begin{boundary}[Exact force of locality]
The theorem says evaluation reads only syntactically named inputs.  It does
not prove that the operations are linear or nonlinear in any analytic sense,
that the tree is stable, that two different trees are equivalent, or that a
solution to a tree equation exists.
\end{boundary}

\subsection{Typed colored extension (roadmap)}

\begin{definition}[Finitary colored operation signature]
\label{def:signature}
A finitary colored signature consists of:
\begin{itemize}
\item a finite type $V$ of colors;
\item a type $\Sigma$ of operation symbols;
\item an arity $r(\sigma)\in\N$ for every $\sigma\in\Sigma$;
\item input colors
$\operatorname{in}_\sigma:\operatorname{Fin}(r(\sigma))\to V$;
\item an output color $\operatorname{out}(\sigma)\in V$.
\end{itemize}
A realization assigns a space $X_v$ to each color and a map
\[
\Phi_\sigma:
\prod_{q\in\operatorname{Fin}(r(\sigma))}X_{\operatorname{in}_\sigma(q)}
\longrightarrow X_{\operatorname{out}(\sigma)}.
\]
\end{definition}
\statusroadmap{The intended Lean target is a dependent inductive typed tree; no such declaration is part of the version 1.0.1 completion contract.}

\begin{definition}[Typed operation tree and evaluation]
\label{def:operation-tree}
A typed operation tree is recursively either a leaf of declared color or a
node labeled by $\sigma$, with one correctly colored child for each input
slot of $\sigma$.  Evaluation is recursive:
\[
\ev(\operatorname{node}(\sigma;t_1,\ldots,t_{r(\sigma)}))
=
\Phi_\sigma(\ev(t_1),\ldots,\ev(t_{r(\sigma)})).
\]
\end{definition}
\statusroadmap{Candidate API: \leanname{Komorebi.PaperIII.NonlinearObstructions.ColoredOperationTree} and \leanname{ColoredOperationTree.eval}.  This is stronger than the compiled uncolored unary--binary grammar.}

\subsection{Packet labelings and the multilinear first target}

Suppose every $X_v$ has a finite complete packet family
$(P^v_i)_{i\in I_v}$.  A packet labeling assigns an admissible packet index
to each edge or node of a typed tree.  For multilinear $\Phi_\sigma$, inserting
\[
\id_{X_v}=\sum_{i\in I_v}P^v_i
\]
at every input and output yields a finite sum over packet-labeled trees.
This is the branching analogue of inserting intermediate projectors in
\cref{prop:heterogeneous-path}.

\begin{roadmap}[Finite multilinear tree expansion]
\label{roadmap:multilinear-tree}
The first milestone for the nonlinear module is an exact theorem saying that the
evaluation of a finite tree of multilinear nodes equals the finite sum of its
packet-resolved labeled evaluations.  Required proof obligations include:
\begin{enumerate}[label=(\alph*)]
\item a dependent well-typed tree syntax;
\item finite packet-label types at every color;
\item an explicit convention for nullary and unary nodes;
\item multilinearity in every input slot;
\item preservation of input-slot order without assuming commutativity of
operations.
\end{enumerate}
\end{roadmap}
\statusroadmap{Candidate target: \leanname{Komorebi.PaperIII.NonlinearObstructions.ColoredOperationTree.eval_eq_sum_packetLabellings}.}

\subsection{Connection defects at nonlinear nodes}

For coordinate models of the source and target colors, an additive target
coordinate space, and a chosen noncanonical connection, a pointwise node
defect can be defined by
\[
\mathfrak D^\Gamma_{b,a}(\Phi)
=
\Gamma^Y_{b\leftarrow a}\circ\widehat\Phi_a
-
\widehat\Phi_b\circ
\left(\prod_q\Gamma^{X_q}_{b\leftarrow a}\right).
\]
This expression compares two \emph{chosen} routes.  With canonical exact
transport in every slot it vanishes by \cref{prop:set-naturality}.

\begin{roadmap}[Operation-tree defect propagation]
\label{roadmap:tree-defect}
For multilinear nodes, one expects an exact finite expansion obtained by
replacing one or more nodes by their defects and transporting the remaining
branches.  For smooth nonlinear nodes, the corresponding chain rule is
indexed by rooted partitions and Fa\`a di Bruno trees.  A formal theorem will
require a fixed differentiability category, base points, derivative
conventions, symmetry factors, and remainder policy.
\end{roadmap}
\statusroadmap{No tree-level defect identity, higher-chain-rule theorem, or nonlinear norm bound is claimed here.}

\begin{boundary}[Analytic and geometric exclusions]
This version does not prove convergence of an infinite tree series, existence
of derivatives, a Banach-space Fa\`a di Bruno theorem, a Baker--Campbell--Hausdorff
expansion, a continuum limit, or an identification of algebraic holonomy with
curvature of a smooth bundle.  It also does not derive a connection from
data, optimization, physics, or probability.
\end{boundary}

\section{Formalization ledger for version 1.0.1}
\label{sec:ledger}

\small
\begin{longtable}{L{0.17\textwidth}L{0.18\textwidth}L{0.25\textwidth}L{0.22\textwidth}}
\toprule
Topic & Status & Declaration or roadmap target & Exact boundary\\
\midrule
\endfirsthead
\toprule
Topic & Status & Declaration or roadmap target & Exact boundary\\
\midrule
\endhead
Packet matrix and edge operator
& Lean-backed III
& \leanname{PacketGraph.packetMatrix}; \leanname{EdgeSystem.edgeOp}
& Finite scalar matrices over a semiring; one scalar coordinate per packet.\\
\addlinespace
Edge support and pair dichotomy
& Lean-backed III
& \leanname{edgeOp_eq_zero_off_support}; \leanname{edgeOp_mul_edgeOp_eq_zero}; \leanname{edgeOp_mul_edgeOp_of_composable}
& Broken interfaces vanish; matched pairs collapse to the endpoint block with coefficient in operator order.\\
\addlinespace
Typed word evaluation
& Lean-backed III
& \leanname{Composable}; \leanname{wordOp}; \leanname{wordOp_append}; \leanname{wordOp_eq_zero_of_not_composable}
& List order is execution order; concatenation reverses operator-product order; nonwalks vanish.\\
\addlinespace
Finite word and walk expansions
& Lean-backed III
& \leanname{explicitWordExpansion}; \leanname{walkExpansion}; \leanname{wordExpansion_eq_explicitWordExpansion}; \leanname{transferOp_pow_eq_walkExpansion}
& Actual finite tuple sum and zero-filtered composable-walk sum; finite graph and packet types.\\
\addlinespace
Endpoint coefficient walk sum
& Lean-backed III
& \leanname{reverseCoefficient}; \leanname{Connects}; \leanname{transferOp_pow_apply_eq_sum_walks}
& Exact scalar endpoint entry, including the empty-walk case and noncommutative coefficient order.\\
\addlinespace
Fixed reachability and orbit support
& Lean-backed III; inherited hull bridge
& \leanname{transferOp_pow_apply_eq_zero_of_no_walk}; \leanname{transferOp_pow_mulVec_apply_eq_zero_of_unreachable_support}
& Finite powers and supported seeds; the orbit-hull span conclusion uses the frozen Paper I contract.\\
\addlinespace
Exact frame transport
& Lean-backed III
& \leanname{exactTransport_cocycle}; \leanname{exactPathTransport_eq_endpoint}; \leanname{exactPathTransport_loop}
& One frame family on a single module; arbitrary finite paths telescope and closed paths are trivial.\\
\addlinespace
Exact coordinate cocycle
& Inherited Lean-backed II
& \leanname{NestedRefinement.refineCoordinates_comp}; \leanname{NestedRefinement.coordinateRefinementEquiv_trans}
& Complete finite projector coordinates, not reduction.\\
\addlinespace
Typed coordinate naturality
& Inherited Lean-backed II
& \leanname{NestedRefinement.coordinateOperator_comp}; \leanname{NestedRefinement.coordinateOperator_naturality}
& Linear maps over semiring modules.\\
\addlinespace
Minimal connection API
& Lean-backed III
& \leanname{Transport.Connection}; \leanname{cocycleDefect}; \leanname{triangleHolonomy}; \leanname{triangleHolonomy_eq_refl_of_flat}
& Only self-normalization is assumed; reverse inversion and the cocycle are not structure fields.\\
\addlinespace
Twisted three-point loop
& Lean-backed III
& \leanname{twistedThreePointConnection}; \leanname{twisted_triangleHolonomy_ne_refl}
& Concrete rational connection; the selected loop sends \(1\) to \(2\).\\
\addlinespace
Reversible additive defect
& Lean-backed III
& \leanname{ReversibleConnection}; \leanname{triangleHolonomy_sub_id}; \leanname{cocycleDefect_eq_zero_iff_triangleHolonomy_eq_refl}
& Heterogeneous module fibers, reverse-edge inversion, and an additive defect over a commutative ring.\\
\addlinespace
Reversible finite paths
& Lean-backed III
& \leanname{pathTransport_eq_direct_of_flat}; \leanname{indexedPathTransport_eq_direct_of_flat}; \leanname{indexedLoopTransport_eq_refl_of_flat}
& Finite pair-connection paths collapse to endpoints under algebraic flatness.\\
\addlinespace
Gauge covariance
& Lean-backed III
& \leanname{gauge_cocycleDefect}; \leanname{gauge_triangleHolonomy}; \leanname{gauge_flat_iff}
& Fiberwise algebraic conjugation; no smooth gauge theory.\\
\addlinespace
Flat versus canonical transport
& Lean-backed III
& \leanname{flatGauge}; \leanname{flat_isGaugeOfCanonical}
& A chosen base chart and ambient frames exhibit every flat pair connection as a gauge transform of canonical transport.\\
\addlinespace
Operation-tree locality
& Lean-backed III
& \leanname{OperationTree}; \leanname{OperationTree.eval}; \leanname{OperationTree.eval_congr_on_inputs}
& Arbitrary unary and binary semantics; only dependence on named inputs is proved.\\
\addlinespace
Generic obstruction packet
& Lean-backed III
& \leanname{ObstructionPacket.not_isSolution}; \leanname{not_isSolution_iff_nonempty_obstructionPacket}
& Exact candidate-specific logical certificate; no global classification.\\
\addlinespace
Holonomy block obstruction
& Lean-backed III
& \leanname{obstruction_sum_blocks}; \leanname{obstructionBlocks_zero_iff}; \leanname{obstructionDetection_equivalences}
& Finite complete idempotent packet family on the base fiber; orthogonality is not required by the compiled reconstruction.\\
\addlinespace
Packet witness and block gauge
& Lean-backed III
& \leanname{exists_packetWitness_of_obstructionBlock_ne_zero}; \leanname{obstructionBlock_gauge_of_commutes}; \leanname{obstructionBlock_gauge_eq_zero_iff}
& Witness extraction is classical; blockwise covariance requires the gauge to commute with every packet projector.\\
\addlinespace
Heterogeneous matrix paths
& Lean-backed III
& \leanname{heterogeneousPacketPathExpansion}; \leanname{heterogeneousProduct_apply_eq_zero_of_no_supportedPath}
& Arbitrarily many different matrices over one common finite packet type; fully typed changing-index chains compile only at lengths two and three.\\
\addlinespace
Arbitrary-function coordinates
& Lean-backed III
& \leanname{exactCoordinateMap_comp}; \leanname{exactCoordinateMap_naturality}
& Exact set-level equivalences and arbitrary functions; no linearity or differentiability assumption.\\
\addlinespace
Colored and smooth trees
& Roadmap only
& roadmap \leanname{ColoredOperationTree.eval_eq_sum_packetLabellings}
& Dependent colors, multilinear packet expansion, derivatives, and Fa\`a di Bruno identities remain outside the contract.\\
\bottomrule
\end{longtable}
\normalsize

\subsection{Kernel-backed Paper III surface}

The authoritative declaration-level contract is
\path{formalization/lean/LEAN_THEOREM_MAP_v1.0.1.md}.  It groups the compiled
surface as follows.
\begin{itemize}
\item \textbf{EdgeOperators and EdgeExtensions:} scalar packet matrices,
typed edge and word support, finite word/walk expansions, endpoint coefficient
sums, fixed-length reachability, orbit support, and common-index
heterogeneous products with declared supports.
\item \textbf{Transport and ReversibleConnections:} exact frame transport,
the minimal connection, heterogeneous reversible pair transport, additive
defect and triangle holonomy, finite paths, canonical connections, and the
fiberwise gauge action.
\item \textbf{CoordinateNaturality and HolonomyBlocks:} composition and
naturality for arbitrary functions between exact coordinate systems, finite
two-sided block reconstruction, equivalent obstruction criteria,
packet-supported witnesses, and packet-preserving block gauges.
\item \textbf{NonlinearObstructions:} finite unary--binary operation trees,
input locality, and candidate-specific obstruction certificates.
\end{itemize}
The release verification report records the checker-generated declaration
and audit totals.  The parity gate requires every fully qualified Paper III
target printed beside a statement to occur in the theorem map, rejects stale
or duplicate targets, and checks it against the compiled environment.
\path{KernelAudit.lean} reports dependencies for the mapped and supporting
theorems, and its fail-closed gate scans every compiled
\leanname{Komorebi.PaperIII} declaration.  The four inherited statement
contracts remain identified separately by their Papers I--II names.

\begin{boundary}[No scope promotion by prose]
The displayed endpoint coefficient formula, reachability results, reversible
connection theorems, additive defect--holonomy identity, gauge covariance,
holonomy-block reconstruction, packet witnesses, common-index heterogeneous
path expansion, and arbitrary-function coordinate naturality are all part of
the version 1.0.1 kernel-backed surface.  What remains outside is the
dependent fiber-valued edge-system generalization, a fully dependent
arbitrary-length chain whose module and packet type change at every stage,
dependent colored operation trees, packet-labeled multilinear tree
expansions, nonlinear tree-defect and Fa\`a di Bruno identities, and the
analytic and geometric extensions listed above.
\end{boundary}

\section{Conclusion}

Typed packet support turns the expansion of a transfer power into a walk
calculus: incompatible words vanish for structural reasons, not because of
numerical cancellation.  Exact packet-coordinate transport is even more
rigid.  It is the canonical analysis--synthesis equivalence and therefore
satisfies its cocycle law exactly; every finite canonical loop telescopes to
the identity.  Holonomy enters only after a separate connection is supplied.
The kernel-backed minimal API proves that flatness kills both the relative
defect and holonomy, and it exhibits a nontrivial rational loop.  Under the
kernel-backed reversible API, a nonzero additive triangle defect is equivalent to nonidentity
holonomy, transforms covariantly under gauge, and is detected exactly by a
finite family of packet blocks.  Finite flat paths collapse to their endpoints,
and every flat reversible pair connection is explicitly gauge-equivalent to
canonical transport through a chosen base chart.

These facts define a clean boundary for the nonlinear program.  At the audited
boundary, exact coordinates preserve composition and satisfy naturality for
arbitrary functions in Paper III, while the typed linear packet-coordinate
counterpart remains inherited from Paper~II.  Nontrivial nonlinear defects
must therefore be attached to additional modeling choices or to
operation-specific packet compatibility.  Colored operation trees,
multilinear packet labelings, and Fa\`a di Bruno expansions remain explicit
future milestones rather than completed results.

\begin{thebibliography}{9}

\bibitem{paperI}
Komorebi P.,
\emph{Projector-Resolved Seed Calculus I: Typed Seeds, Operator Grammars, and
Exact Reconstruction},
version 1.2.1, 2026.

\bibitem{paperII}
Komorebi P.,
\emph{Projector-Resolved Seed Calculus II: Polyphase Refinement, Weighted
Transfer, and Dyadic--Odd Resolution},
version 1.2.1, 2026.

\bibitem{lean}
Leonardo de Moura and Sebastian Ullrich,
\emph{The Lean 4 Theorem Prover and Programming Language},
in Andr\'e Platzer and Geoff Sutcliffe (eds.),
\emph{Automated Deduction -- CADE 28}, Lecture Notes in Computer Science,
vol.~12699, pp.~625--635, Springer, Cham, 2021.
DOI: \href{https://doi.org/10.1007/978-3-030-79876-5_37}%
{\texttt{10.1007/978-3-030-79876-5\_37}}.

\bibitem{mathlib}
The mathlib Community,
\emph{The Lean mathematical library}, in
\emph{Proceedings of the 9th ACM SIGPLAN International Conference on
Certified Programs and Proofs (CPP 2020)}, pp.~367--381,
Association for Computing Machinery, New York, NY, USA, 2020.
DOI: \href{https://doi.org/10.1145/3372885.3373824}%
{\texttt{10.1145/3372885.3373824}}.

\end{thebibliography}

\end{document}
